\section{Incidence hexade–octade et fonctionnelle orientée}
\subsection{Incidence dodécaphasique et fonctionnelle de retour}

La contribution de Barbero--Immirzi utilise le transport angulaire et
l’incidence hexade--octade. Soient $\mathsf H\subset\mathsf O$ une hexade
et une octade avec la paire marquée de l’hexade conservée. Les ensembles de
drapeaux pertinents sont
\begin{equation}
 \mathcal F_6=\mathsf H\times\binom{\mathsf H}{2}
 \lhook\joinrel\longrightarrow
 \mathcal F_8=\mathsf O\times\binom{\mathsf H}{2},
 \qquad |\mathcal F_6|=90,\quad |\mathcal F_8|=120.
 \label{iii:eq:barbero-flags}
\end{equation}
Les cardinaux sont $6\binom62$ et $8\binom62$ ; la même incidence
conserve le défaut normalisé
$(90-120)/(24\binom62)=-1/12$. Ce dénombrement identifie les deux secteurs
employés par le retour, tandis que le calendrier détermine leurs
multiplicités relatives.

Pour $m\in\{90,120\}$, le retour tritique sur un canal contractant est
\begin{equation}
 S_m(q)=\sum_{j\geq0}q^{m+3mj}
       =\frac{q^m}{1-q^{3m}},\qquad 0<q<1.
 \label{iii:eq:barbero-return-series}
\end{equation}
Chaque terme correspond à la hauteur initiale $m$ et à $j$ retours de
longueur $3m$. La série géométrique conserve donc tant la hauteur
initiale que la période du retour. Le calendrier dodécaphasique
$\mathbb Z/12\mathbb Z$ apporte douze secteurs et une unique incidence
orientée de frontière, notée $\partial_{\rm ret}$. Dans le module
libre des événements, sa chaîne fondamentale est
\begin{equation}
 c_{12}^{+}
 =\sum_{j\in\mathbb Z/12\mathbb Z}[j,\mathcal F_6]
  -[\partial_{\rm ret},\mathcal F_8].
 \label{iii:eq:barbero-cycle}
\end{equation}
Le lecteur linéaire $\mathscr L_q$ assigne $S_{90}(q)$ à chaque générateur
$[j,\mathcal F_6]$ et $S_{120}(q)$ au générateur de frontière. On obtient
\begin{equation}
 \mathscr L_q(c_{12}^{+})=12S_{90}(q)-S_{120}(q).
 \label{iii:eq:barbero-weights}
\end{equation}
Les poids $12$ et $-1$ sont les multiplicités orientées de la chaîne.
L’orientation conjuguée satisfait $c_{12}^{-}=-c_{12}^{+}$, avec les
mêmes types d’incidence.

\begin{proposition}[Coefficient singulier du retour fondamental]
\label{iii:prop:barbero-pole}
Le coefficient de $(1-q)^{-1}$ de
$\mathscr L_q(c_{12}^{+})$ est $1/24$.
\end{proposition}
\begin{proof}
Pour chaque entier $m>0$, la factorisation
$1-q^{3m}=(1-q)\sum_{j=0}^{3m-1}q^j$ donne
\[
 \lim_{q\uparrow1}(1-q)S_m(q)=\frac1{3m}.
\]
Appliquer cette identité à la fonctionnelle déjà déterminée dans
\eqref{iii:eq:barbero-weights} produit
\begin{equation}
 \lim_{q\uparrow1}(1-q)\mathscr L_q(c_{12}^{+})
 =\frac{12}{270}-\frac1{360}=\frac1{24}.
 \label{iii:eq:barbero-pole}
\end{equation}
Dans la coordonnée complexe $q-1$, le résidu est $-1/24$.
\end{proof}

L’évaluation singulière vérifie une conséquence de l’incidence et du
retour tritique. L’ordre démonstratif est fixé par
\eqref{iii:eq:barbero-flags}--\eqref{iii:eq:barbero-weights} ; le
coefficient singulier se calcule après la construction de la chaîne.

